\documentclass[11pt,a4paper]{article}
\usepackage[margin=2.4cm]{geometry}
\usepackage{amsmath,amssymb,amsthm,mathtools}
\usepackage{booktabs}
\usepackage{enumitem}
\usepackage{xcolor}
\usepackage{hyperref}

\hypersetup{
  colorlinks=true,
  linkcolor=blue!70!black,
  citecolor=blue!70!black,
  urlcolor=blue!70!black,
  pdftitle={Cuboctahedral Rotator Calculus, Part I: Discrete Euler Exponential on Z_512, Universal Constant Wronskian, and Closed-Form Quadrature on the Cuboctahedron O_2(3) in Lean 4},
  pdfauthor={Dmitrii A. Karataev}
}

\theoremstyle{plain}
\newtheorem{theorem}{Theorem}[section]
\newtheorem{lemma}[theorem]{Lemma}
\newtheorem{proposition}[theorem]{Proposition}
\newtheorem{corollary}[theorem]{Corollary}
\theoremstyle{definition}
\newtheorem{definition}[theorem]{Definition}
\newtheorem{remark}[theorem]{Remark}

\title{\textbf{Cuboctahedral Rotator Calculus, Part~I:\\ Discrete Euler Exponential on $\mathbb{Z}_{512}$, Universal Constant Wronskian, and Closed-Form Quadrature on the Cuboctahedron $O_2(3)$ in Lean~4}}
\author{Dmitrii A.~Karataev}
\date{}

\begin{document}
\maketitle

\begin{abstract}
We investigate discrete trigonometric and integral calculus on the three-dimensional cuboctahedral shell $O_2(d)\subset\{-1,0,+1\}^d$ ($d=3$, $|V_{O_2}|=12$ vertices, $|\square|=6$ square faces, $|\triangle|=8$ triangular faces) and the associated finite phase ring $\mathbb{Z}_{N_{\mathrm{ring}}}$ ($N_{\mathrm{ring}} = 4\,Q_\square = 512$, where $Q_\square = 2^{|\triangle|-1} = 128$ is the quarter-period). Extending the piecewise-linear calculus of Thompson (\texttt{arXiv:1205.3814}), where piecewise derivatives lost the classical sine--cosine coupling at diamond corners and integral calculus remained open, we prove that the branchless projection pair $\mathcal{E}_d(\varphi) = (\mathcal{C}_d(\varphi),\mathcal{S}_d(\varphi))$ forms an exact discrete Euler exponential on the $512$-gon generating the five fundamental constants $\{0,1,i,\pi,e\}$: a quarter-period shift $+Q_\square(d)=+128$ acts across all $512$ ring points as exact multiplication by the imaginary unit~$i$, a half-period shift $+\Phi_{\mathrm{half}}(d)=+256$ satisfies Euler's identity $\mathcal{E}_d(\varphi+\Phi_{\mathrm{half}}(d))+\mathcal{E}_d(\varphi)=(0,0)$ ($e^{i\pi}+1=0$), and the one-step difference Wronskian $\mathcal{W}_d(\varphi)=\mathcal{C}_d(\varphi)\mathcal{S}_d(\varphi+1)-\mathcal{S}_d(\varphi)\mathcal{C}_d(\varphi+1)=\langle \mathcal{E}_d(\varphi+Q_\square(d)),\Delta\mathcal{E}_d(\varphi)\rangle$ is identically equal to the constant $Q_\square(d)=128$ at all $512$ points of the ring, including the four corner vertices.

We further prove that the two-tick symplectic Verlet difference recurrence $\varphi_{t+1}-2\varphi_t+\varphi_{t-1}=\Phi$ unifies Liouville's elementary circular basis ($\Phi=0$) with the non-elementary quadratic-phase Fresnel--Airy basis ($\Phi\ne 0$), while two-ring coupling at the intrinsic cuboctahedron modulus $k^2 = 1/\|e\|_2^2 = 1/(d-1) = 1/2$ generates Gauss's lemniscate Jacobi elliptic triad $(\mathrm{sn}_d,\mathrm{cn}_d,\mathrm{dn}_d)$. This structure yields exact closed-form integer evaluations of the discrete Fresnel diffraction integrals ($\mathcal{C}_{\mathrm{Fresnel}}=\mathcal{S}_{\mathrm{Fresnel}}=2048$), complete and incomplete elliptic integrals of the second kind, the Dirichlet sine integral ($\mathrm{Si}_3(128)=128$), the Parseval energy integral ($2,796,288$), and the zero-order Bessel integral ($\mathcal{J}_{0,3}(0)=65,536$, $\mathcal{J}_{0,3}(1)=32,768$, and exact root $\mathcal{J}_{0,3}(4)=0$). Finally, projecting the $12$ vertices of $O_2(3)$ onto any coordinate axis yields the Archimedes equizonal partition $(4,4,4)$, which together with the central node of weight $12$ derives Simpson's parabolic weights $(4,16,4)=4\times(1,4,1)$ and the exact 3D continuum Riemann integral $\iiint_{[-1,1]^3} F\,\mathrm{d}V$. All definitions and theorems (\texttt{thm1}, \texttt{thm2}, \texttt{thm3}) are formally verified in the \textbf{Lean~4} theorem prover (\texttt{0 sorryAx}, \texttt{0 Classical.choice}).

\medskip\noindent
\textbf{Keywords:} piecewise-linear trigonometry, Euler's identity, difference Wronskian, Weierstrass substitution, cuboctahedron, Fresnel integrals, lemniscate Jacobi elliptic functions, formal verification, Lean~4.

\smallskip\noindent
\textbf{MSC 2020:} 26A06, 33E05, 39A12, 65D30, 68V20.
\end{abstract}

\section{Introduction}
\label{sec:intro}

Piecewise-linear trigonometry in the $L^1$ metric (``taxicab trigonometry'' on the diamond $|x|+|y|=R$) was developed by Krause~\cite{Krause1986} and Thompson~\cite{Thompson2011,Thompson2012}. In~\cite{Thompson2012} (\texttt{arXiv:1205.3814}), Thompson investigated the piecewise derivatives of the $L^1$ sine, cosine, tangent, and secant functions. However, three structural obstacles remained unresolved in~\cite{Thompson2012}:
\begin{enumerate}[leftmargin=*,noitemsep]
\item the piecewise-constant derivatives of the $L^1$ sine and cosine are undefined at the four corner extrema and lose the classical Euclidean differential coupling $(\sin\theta)'=\cos\theta$, $(\cos\theta)'=-\sin\theta$;
\item representing the piecewise branches across quadrants required complex powers $i^{k-2}$ branched by the parity of the quadrant index $k$;
\item integral calculus, the lift to the Euclidean ($L^2$) circle, and the evaluation of classical non-elementary integrals (Fresnel, Jacobi elliptic, Bessel, Dirichlet) were left open.
\end{enumerate}

In this paper (Part~I of a four-part series), we resolve these obstacles by formulating the rotator over the three-dimensional Archimedean cuboctahedral shell $O_2(d) = \{e=(e_x,e_y,e_z)\in\{-1,0,+1\}^3 \mid e_x^2+e_y^2+e_z^2 = d-1\}$ ($d=3$) and the associated modular integer phase ring $\mathbb{Z}_{N_{\mathrm{ring}}(d)}$:
\[
  |{\triangle}(d)| = 2^d = 8,\qquad
  |\square(d)| = 2d = 6,\qquad
  |V_{O_2}(d)| = 2d(d-1) = 12,
\]
\[
  d_{\mathrm{in}}(\square) = 4,\qquad
  d_{\mathrm{out}}(E)(d) = 2^d - d = 5,\qquad
  \|e\|_2^2 = d - 1 = 2,
\]
\[
  Q_\square(d) = 2^{|{\triangle}(d)|-1} = 128,\qquad
  \Phi_{\mathrm{half}}(d) = 2\,Q_\square(d) = 256,\qquad
  N_{\mathrm{ring}}(d) = 4\,Q_\square(d) = 512,\qquad
  K_P(d) = \frac{N_{\mathrm{ring}}(d)}{|{\triangle}(d)|} = 64.
\]
All definitions and theorems below correspond to Theorems~1--3 (\texttt{thm1}--\texttt{thm3}) in the self-contained \textbf{Lean~4}~\cite{Moura2021} module \texttt{RotatorCalculusPaper.lean}.

\begin{definition}[Model-Theoretic Two-Sorted Algebraic System $\langle\mathcal{M},\mathcal{T}\rangle$ in the Sense of Mal'cev~\cite{Malcev1973}]
\label{def:malcev-carriers}
Throughout this series, the term \emph{carrier} (\emph{носитель} in the Russian model-theoretic tradition of A.\,I.~Mal'cev, \emph{Algebraic Systems}~\cite{Malcev1973}) is used strictly in its standard model-theoretic meaning: the underlying domain set of an algebraic structure $\mathfrak{A} = \langle A;\,\Sigma\rangle$. We formulate the calculus as a two-sorted algebraic structure $\mathfrak{S} = \langle \mathcal{M},\,\mathcal{T};\,\pi_{M\to T},\,\iota_{T\to M}\rangle$ with two carriers:
\begin{enumerate}[leftmargin=*,noitemsep]
\item the \textbf{discrete $M$-carrier} $\mathcal{M} = \mathbb{Z}_{512}$ (equipped with $L^1$ diamond projections $\mathcal{C}_3,\mathcal{S}_3$ and modular arithmetic over the cuboctahedral lattice $\mathbb{Z}^3$);
\item the \textbf{continuous (Euclidean) $T$-carrier} $\mathcal{T}$ (the continuum sort of $L^2$ Euclidean circle invariants $X^2+Y^2=R^2$, rational/real projective coordinates $\mathbb{Q}\subset\mathbb{R}$, and continuum Riemann integrals);
\item connected by the root-free Weierstrass lift $\pi_{M\to T}(C,S) = (C^2-S^2,2CS,C^2+S^2)$, linear retraction $\iota_{T\to M}$, and stereographic homeomorphism $[0,512]/512 = [0,1] \cong \mathbb{R}$.
\end{enumerate}
\end{definition}

\section{Discrete Euler Exponential, Euler's Identity \texorpdfstring{$e^{i\pi}+1=0$}{e\^{}(i pi)+1=0}, Constant Wronskian, and the Weierstrass 512-gon Lift}
\label{sec:wronskian-euler}

For any $x\in\mathbb{Z}$, define the principal symmetric branch wrapper into $(-\Phi_{\mathrm{half}}(d),\Phi_{\mathrm{half}}(d)]$:
\begin{equation}
  \label{eq:wrap-ring}
  \mathrm{wrap}_d(x) \;:=\; \Phi_{\mathrm{half}}(d) - \bigl((\Phi_{\mathrm{half}}(d) - x) \bmod N_{\mathrm{ring}}(d)\bigr),
\end{equation}
and the branchless piecewise-linear cosine and sine projections:
\begin{equation}
  \label{eq:proj-cs}
  \mathcal{C}_d(\varphi) \;:=\; Q_\square(d) - |\mathrm{wrap}_d(\varphi)|,\qquad
  \mathcal{S}_d(\varphi) \;:=\; \mathcal{C}_d\bigl(\varphi - Q_\square(d)\bigr).
\end{equation}

\begin{definition}[Discrete Euler exponential and one-step difference Wronskian]
\label{def:euler-exp-wronskian}
We define the \emph{discrete Euler exponential} on $\mathbb{Z}_{N_{\mathrm{ring}}(d)}$ by the pair
\[
  \mathcal{E}_d(\varphi) \;:=\; \bigl(\mathcal{C}_d(\varphi),\,\mathcal{S}_d(\varphi)\bigr) \;\in\; \mathbb{Z}^2,
\]
and the \emph{one-step difference Wronskian} by
\begin{equation}
  \label{eq:rotator-wronskian}
  \mathcal{W}_d(\varphi) \;:=\; \mathcal{C}_d(\varphi)\,\mathcal{S}_d(\varphi+1) - \mathcal{S}_d(\varphi)\,\mathcal{C}_d(\varphi+1)
  \;=\; \mathcal{C}_d(\varphi)\,\Delta\mathcal{S}_d(\varphi) - \mathcal{S}_d(\varphi)\,\Delta\mathcal{C}_d(\varphi),
\end{equation}
where $\Delta f(\varphi) = f(\varphi+1)-f(\varphi)$.
\end{definition}

To each $\varphi\in\mathbb{Z}_{N_{\mathrm{ring}}(d)}$ we associate the rational Brahmagupta--Weierstrass lift onto the Euclidean circle $X^2+Y^2=R^2$:
\begin{equation}
  \label{eq:weierstrass-lift}
  X_d(\varphi) := \mathcal{C}_d(\varphi)^2 - \mathcal{S}_d(\varphi)^2,\qquad
  Y_d(\varphi) := \|e\|_2^2\,\mathcal{C}_d(\varphi)\,\mathcal{S}_d(\varphi),\qquad
  R_d(\varphi) := \mathcal{C}_d(\varphi)^2 + \mathcal{S}_d(\varphi)^2.
\end{equation}

\begin{theorem}[Discrete Euler's Identity, Universal Constant Wronskian, Weierstrass Lift, and Fundamental Theorem of Discrete Calculus]
\label{thm:main-1}
For $d=3$, the following exact identities hold across all $N_{\mathrm{ring}}(d)=512$ points $\varphi\in\mathbb{Z}_{N_{\mathrm{ring}}(d)}$ (\texttt{thm1\_constant\_wronskian\_weierstrass\_lift\_and\_newton\_leibniz}):
\begin{enumerate}[label=\textbf{(\roman*)},leftmargin=*]
\item \textbf{Generation of $\{0,1,i,\pi,e\}$ and Euler's identity $e^{i\pi}+1=0$:}
  \begin{itemize}[noitemsep]
  \item the $L^1$ norm is constant: $|\mathcal{C}_d(\varphi)| + |\mathcal{S}_d(\varphi)| \equiv Q_\square(d)$, with $\mathcal{E}_d(0) = (Q_\square(d),0)$ representing unity~$1$;
  \item a quarter-period shift $+Q_\square(d)$ acts across all $512$ points as exact multiplication by the imaginary unit~$i$:
    \[
      \mathcal{E}_d\bigl(\varphi + Q_\square(d)\bigr) \;=\; \bigl(-\mathcal{S}_d(\varphi),\; +\mathcal{C}_d(\varphi)\bigr) \;=\; i\cdot\mathcal{E}_d(\varphi);
    \]
  \item a half-period shift $+2Q_\square(d)=+\Phi_{\mathrm{half}}(d)$ ($\pi$) satisfies Euler's identity $e^{i(\varphi+\pi)}+e^{i\varphi}=0$ identically across the ring:
    \[
      \mathcal{E}_d\bigl(\varphi + \Phi_{\mathrm{half}}(d)\bigr) + \mathcal{E}_d(\varphi) \;=\; (0,0),
      \qquad
      \mathcal{E}_d\bigl(\Phi_{\mathrm{half}}(d)\bigr) + \mathcal{E}_d(0) \;=\; (0,0);
    \]
  \item the inner product of $i\mathcal{E}_d(\varphi)=\mathcal{E}_d(\varphi+Q_\square(d))$ with the difference derivative $\Delta\mathcal{E}_d(\varphi)$ coincides with the Wronskian $\mathcal{W}_d(\varphi)$ and is \textbf{identically equal to the constant $Q_\square(d)$ at all $512$ points of the ring (including the four corner vertices)}:
    \begin{equation}
      \label{eq:constant-wronskian-all-512}
      \bigl\langle \mathcal{E}_d\bigl(\varphi+Q_\square(d)\bigr),\,\Delta\mathcal{E}_d(\varphi)\bigr\rangle
      \;=\; \mathcal{W}_d(\varphi) \;\equiv\; Q_\square(d) \quad (=128),\qquad
      \sum_{\varphi=0}^{N_{\mathrm{ring}}(d)-1}\mathcal{W}_d(\varphi) = N_{\mathrm{ring}}(d)\,Q_\square(d).
    \end{equation}
  \end{itemize}
\item \textbf{Root-free Euclidean Weierstrass $512$-gon and arc-step differential:}
  at all $N_{\mathrm{ring}}(d)$ points, $X_d(\varphi)^2 + Y_d(\varphi)^2 = R_d(\varphi)^2$, and the one-step Euclidean cross and dot products are governed by $\mathcal{I}_d(\varphi) := \mathcal{C}_d(\varphi)\mathcal{C}_d(\varphi+1) + \mathcal{S}_d(\varphi)\mathcal{S}_d(\varphi+1)$:
  \begin{equation}
    \label{eq:brahmagupta-step}
    X_d(\varphi)Y_d(\varphi+1) - Y_d(\varphi)X_d(\varphi+1) \;=\; \|e\|_2^2\,Q_\square(d)\,\mathcal{I}_d(\varphi),\qquad
    X_d(\varphi)X_d(\varphi+1) + Y_d(\varphi)Y_d(\varphi+1) \;=\; \mathcal{I}_d(\varphi)^2 - Q_\square(d)^2.
  \end{equation}
\item \textbf{Fundamental Theorem of Discrete Calculus (Newton--Leibniz):}
  the piecewise-quadratic antiderivatives $\widehat{\Phi}_{\mathcal{C}}(d,\Phi) := \sum_{\varphi=0}^{\Phi-1}\mathcal{C}_d(\varphi)$ and $\widehat{\Phi}_{\mathcal{S}}(d,\Phi) := \sum_{\varphi=0}^{\Phi-1}\mathcal{S}_d(\varphi)$ satisfy $\Delta\widehat{\Phi}_{\mathcal{C}}(d,\varphi) = \mathcal{C}_d(\varphi)$, $\Delta\widehat{\Phi}_{\mathcal{S}}(d,\varphi) = \mathcal{S}_d(\varphi)$, vanish over the full period $N_{\mathrm{ring}}(d)$, and yield the exact half-wave area $\widehat{\Phi}_{\mathcal{S}}(d,\Phi_{\mathrm{half}}(d)) = Q_\square(d)^2 = 16,384$.
\end{enumerate}
\end{theorem}

\begin{proof}
In the first quadrant $\varphi\in[0,Q_\square(d))$, $\mathcal{C}_d(\varphi)=Q_\square(d)-\varphi$ and $\mathcal{S}_d(\varphi)=\varphi$, so
\[
  \mathcal{W}_d(\varphi) = (Q_\square(d)-\varphi)(\varphi+1) - \varphi(Q_\square(d)-\varphi-1) = Q_\square(d),
\]
as the quadratic terms $-\varphi^2+\varphi^2$ and linear terms cancel identically. By determinant invariance under the quarter-turn rotation $(x,y)\mapsto(-y,x)$, $\mathcal{W}_d(\varphi)\equiv Q_\square(d)$ holds across all four quadrants. Dividing by $\mathcal{C}_d(\varphi)\mathcal{C}_d(\varphi+1)$ yields the exact discrete tangent difference $\Delta\mathcal{T}_d(\varphi) = Q_\square(d)/(\mathcal{C}_d(\varphi)\mathcal{C}_d(\varphi+1))$, turning rational trigonometric sums into telescoping Newton--Leibniz sums. The complete formal proof is checked by kernel reduction (\texttt{rfl}) in Lean~4.
\end{proof}

\section{Native Non-Elementary Basis and Closed-Form Evaluation of Fresnel, Jacobi, Dirichlet, Parseval, and Bessel Integrals}
\label{sec:integrals}

In Liouville's differential algebra, Fresnel integrals $\int\cos(\varphi^2)\,\mathrm{d}\varphi$, elliptic integrals, Dirichlet's sine integral $\mathrm{Si}(\Phi)$, and Bessel integrals $J_0(k)$ admit no closed form in terms of a single circular rotator. On $\mathbb{Z}_{N_{\mathrm{ring}}(d)}$, two native operators close these families:
\begin{enumerate}[leftmargin=*,noitemsep]
\item The two-tick symplectic Verlet difference recurrence
  \begin{equation}
    \label{eq:verlet-parabola}
    \varphi(t+1) - 2\varphi(t) + \varphi(t-1) = \Phi_{\mathrm{accel}}
    \quad\Longrightarrow\quad
    \varphi(t) = \varphi_0 + \omega_0 t + \Phi_{\mathrm{accel}}\,\frac{t(t-1)}{2}
  \end{equation}
  generates Liouville's elementary circular basis at zero acceleration ($\Phi_{\mathrm{accel}}=0$) and the quadratic-phase Fresnel--Airy basis at constant non-zero acceleration ($\Phi_{\mathrm{accel}}\ne 0$).
\item On two coupled rings at the intrinsic cuboctahedron modulus $k^2 = 1/\|e\|_2^2 = 1/(d-1) = 1/2$ (Gauss's lemniscate modulus with complex multiplication in $\mathbb{Z}[i]$), the Jacobi triad $\mathrm{sn}_d(\varphi):=\mathcal{S}_d(\varphi)$, $\mathrm{cn}_d(\varphi):=\mathcal{C}_d(\varphi)$, $2\,\mathrm{dn}_d(\varphi):=\Phi_{\mathrm{half}}(d)-|\mathcal{S}_d(\varphi)|$ satisfies both conservation laws in exact integers.
\end{enumerate}

Define the five canonical ring integrals:
\begin{align}
  \mathcal{C}_{\mathrm{Fresnel}}(d) &:= \sum_{\varphi=0}^{N_{\mathrm{ring}}(d)-1}\mathcal{C}_d(\varphi^2),\qquad
  \mathcal{S}_{\mathrm{Fresnel}}(d) := \sum_{\varphi=0}^{N_{\mathrm{ring}}(d)-1}\mathcal{S}_d(\varphi^2), \label{eq:def-fresnel}\\
  \mathcal{I}_{\mathrm{ell}}(d;a,b) &:= \sum_{\varphi=0}^{N_{\mathrm{ring}}(d)-1}\bigl(a\,|\mathcal{C}_d(\varphi)| + b\,|\mathcal{S}_d(\varphi)|\bigr), \label{eq:def-elliptic}\\
  \mathrm{Si}_d(\Phi) &:= \sum_{\varphi=1}^{\Phi}\mathrm{tdiv}\bigl(\mathcal{S}_d(\varphi),\,\varphi\bigr), \label{eq:def-dirichlet}\\
  \mathcal{I}_{\mathrm{Parseval}}(d) &:= \sum_{\varphi=0}^{N_{\mathrm{ring}}(d)-1}\mathcal{C}_d(\varphi)^2 = \sum_{\varphi=0}^{N_{\mathrm{ring}}(d)-1}\mathcal{S}_d(\varphi)^2, \label{eq:def-parseval}\\
  \mathcal{J}_{0,d}(k) &:= \sum_{\varphi=0}^{N_{\mathrm{ring}}(d)-1}\mathcal{C}_d\bigl(k\,\mathcal{S}_d(\varphi)\bigr). \label{eq:def-bessel}
\end{align}

\begin{theorem}[Closed-Form Fresnel, Jacobi Elliptic, Dirichlet, Parseval, and Bessel Integrals]
\label{thm:main-2}
For $d=3$, the following exact integral evaluations hold on $\mathbb{Z}_{N_{\mathrm{ring}}(d)}$ (\texttt{thm2\_closed\_form\_fresnel\_jacobi\_dirichlet\_parseval\_bessel\_integrals}):
\begin{enumerate}[label=\textbf{(\roman*)},leftmargin=*]
\item \textbf{Lemniscate Jacobi elliptic triad and elliptic integrals:}
  across all $N_{\mathrm{ring}}(d)$ points, $|\mathrm{sn}_d(\varphi)|+|\mathrm{cn}_d(\varphi)|=Q_\square(d)$ and $2\,\mathrm{dn}_d(\varphi)+|\mathrm{sn}_d(\varphi)|=\Phi_{\mathrm{half}}(d)$, with complete elliptic integral $\mathcal{I}_{\mathrm{ell}}(d;\|e\|_2^2,d) = \Phi_{\mathrm{half}}(d)\,Q_\square(d)\,d_{\mathrm{out}}(E)(d) = 163,840$ and first-quadrant incomplete integral $Q_\square(d)^2 = 16,384$.
\item \textbf{Exact stationary-phase evaluation of the Fresnel diffraction integrals:}
  both complete Fresnel integrals coincide and equal
  \begin{equation}
    \label{eq:fresnel-2048}
    \mathcal{C}_{\mathrm{Fresnel}}(d) \;=\; \mathcal{S}_{\mathrm{Fresnel}}(d) \;=\; d_{\mathrm{in}}(\square)\cdot N_{\mathrm{ring}}(d) \;=\; 2048,
  \end{equation}
  concentrated at the stationary-phase bisector $\psi = K_P(d) = 64$ ($\pi/4$), where $\mathcal{C}_d(K_P(d)) = \mathcal{S}_d(K_P(d)) = K_P(d)$.
\item \textbf{Dirichlet sine integral and Parseval energy integral:}
  \[
    \mathrm{Si}_d\bigl(Q_\square(d)\bigr) = \mathrm{Si}_d\bigl(\Phi_{\mathrm{half}}(d)\bigr) = Q_\square(d) = 128,
  \]
  \[
    |\square(d)|\cdot\mathcal{I}_{\mathrm{Parseval}}(d) \;=\; N_{\mathrm{ring}}(d)\,\bigl(\|e\|_2^2\,Q_\square(d)^2 + 1\bigr)
    \quad\Longrightarrow\quad
    \mathcal{I}_{\mathrm{Parseval}}(3) = \frac{512\times 32,769}{6} = 2,796,288,
  \]
  where $2\cdot 128^2+1 = 32,769 = 3\times 10,923$ is divisible by $d=3$.
\item \textbf{Discrete Bessel integral $\mathcal{J}_{0,d}(k)$ and its exact root:}
  \[
    \mathcal{J}_{0,d}(0) = N_{\mathrm{ring}}(d)\,Q_\square(d) = 65,536,\qquad
    \mathcal{J}_{0,d}(1) = \Phi_{\mathrm{half}}(d)\,Q_\square(d) = 32,768,\qquad
    \mathcal{J}_{0,d}\bigl(d_{\mathrm{in}}(\square)\bigr) = \mathcal{J}_{0,3}(4) = 0.
  \]
\end{enumerate}
\end{theorem}

\begin{proof}
Partition $\varphi\in\{0,\dots,511\}$ by $2$-adic valuation: for the $256$ odd $\varphi=2k+1$, the $128$ values $\varphi=2(2k+1)$, and the $64$ values $\varphi=4(2k+1)$, the quadratic images $\varphi^2\bmod 512$ are invariant under the half-period shift $+256$ and cancel to zero by Euler's identity $\mathcal{E}_3(\psi+256)=-\mathcal{E}_3(\psi)$; for the $32$ multiples $\varphi=16k$, $16$ land at $\psi=0$ ($+128$) and $16$ at $\psi=256$ ($-128$), summing to zero. Only the $32$ elements $\varphi=8(2k+1)$ survive, and for every one of them $(2k+1)^2\equiv 1\pmod 8$ implies $\varphi^2\equiv 64 = K_P(3)\pmod{512}$, giving $32\times 64 = 2048$.
\end{proof}

\section{Emergence of the Continuum Circle, \texorpdfstring{$\pi$}{pi}, and the Riemann Integral from 512-gon Rectification and the Archimedes Equizonal Partition of the Cuboctahedron}
\label{sec:continuum-emergence}

For any 3D function $F\colon\mathbb{Z}^3\to\mathbb{Z}$, define the natural center-plus-shell integrator on the cuboctahedron ball $B_1(d)=\{0\}\cup O_2(d)$, weighting the origin $(0,0,0)$ by $|V_{O_2}(d)|=12$ equally against the sum over the $12$ vertices of $O_2(d)$:
\begin{equation}
  \label{eq:cuboct-3d-integrator}
  \widehat{\mathcal{I}}_{O_2(d)}(F) \;:=\; |V_{O_2}(d)|\cdot F(0,0,0) \;+\; \sum_{e\in O_2(d)} F(e_x,e_y,e_z).
\end{equation}

\begin{theorem}[Geometric Derivation of the Continuum Riemann Integral and Simpson's Weights $(1,4,1)/6$ from the Archimedes Equizonal Partition of the Cuboctahedron]
\label{thm:main-3}
For $d=3$, the following relations hold on the cuboctahedral shell $O_2(d)$ and phase ring $\mathbb{Z}_{N_{\mathrm{ring}}(d)}$ (\texttt{thm3\_continuum\_integration\_from\_cuboctahedron\_archimedes\_layers}):
\begin{enumerate}[label=\textbf{(\roman*)},leftmargin=*]
\item \textbf{Archimedes equizonal partition of $O_2(d)$:}
  slicing the $12$ vertices of $O_2(d)=\mathrm{cubicOrbit3D}(d-1)$ along any coordinate axis $e_x\in\{-1,0,+1\}$ partitions the shell into three equal zones of $d_{\mathrm{in}}(\square)=4$ vertices each:
  \[
    \mathrm{cuboctLayerCount}(d,-1) = \mathrm{cuboctLayerCount}(d,0) = \mathrm{cuboctLayerCount}(d,+1) = d_{\mathrm{in}}(\square) = 4.
  \]
  Combining the equatorial zone $e_x=0$ ($4$ vertices) with the central node weight $|V_{O_2}(d)|=12$ in~\eqref{eq:cuboct-3d-integrator} by direct addition 4 + 12 = 16 assembles the slice weights $(-1:0:+1)$ into $(4,\,16,\,4) = 4\times(1,4,1)$ with total weight $\|e\|_2^2\,|V_{O_2}(d)| = d_{\mathrm{in}}(\square)\,|\square(d)| = 24 = 4\times 6$.
\item \textbf{Exact agreement with the 1D and 3D continuum Riemann integral $\iiint_{[-1,1]^d} F(x,y,z)\,\mathrm{d}V$:}
  for all degrees $0\le p\le d$ and 3D quadratic/cubic forms, the cuboctahedron operator~\eqref{eq:cuboct-3d-integrator} reproduces the exact continuum integral:
  \[
    \frac{1}{\|e\|_2^2\,|V_{O_2}(d)|}\,\widehat{\mathcal{I}}_{O_2(d)}(x^p) \;=\; \frac{1}{2}\int_{-1}^{+1} x^p\,\mathrm{d}x \quad (0\le p\le 3),
    \qquad
    \frac{1}{\|e\|_2^2\,|V_{O_2}(d)|}\,\widehat{\mathcal{I}}_{O_2(d)}(x^2+y^2+z^2) \;=\; \frac{1}{2^3}\iiint_{[-1,1]^3}(x^2+y^2+z^2)\,\mathrm{d}V = 1.
  \]
\item \textbf{Remainder-free Pochhammer falling-factorial integration and $512$-gon radius bound:}
  for falling factorials $m^{\underline{k}} = m(m-1)\cdots(m-k+1)$, $d_{\mathrm{in}}(\square)\sum_{m=0}^{|V_{O_2}(d)|} m^{\underline{d}} = (|V_{O_2}(d)|+1)^{\underline{d_{\mathrm{in}}(\square)}}$, the Weierstrass radius satisfies $\|e\|_2^2\,R_d(\varphi) \ge Q_\square(d)^2$ at all $512$ points, and the $512$-gon perimeter defect denominator equals $|\square(d)|\,N_{\mathrm{ring}}(d)^2 = 1,572,864$.
\item \textbf{$\pi(n)$ as a monotone rational function of polygon side count $n$, the Wronskian-cube identity, and calculator-free rational trigonometry:}
  on the Weierstrass $n$-gon with step $s = N_{\mathrm{ring}}(d)/n$, the inscribed and circumscribed rational estimators $\pi_{\mathrm{in}}(n)$ and $\pi_{\mathrm{out}}(n)$ obey the exact \emph{Wronskian-cube bracket identity}
  \begin{equation}
    \label{eq:wronskian-cube-pi-gap}
    \pi_{\mathrm{out}}(n) - \pi_{\mathrm{in}}(n) \;=\; \sum_{j=0}^{n/2-1}\frac{W_s(js)^3}{I_s(js)\,R_d(js)\,R_d((j+1)s)} \;>\; 0,
  \end{equation}
  so under every polygon doubling $n\mapsto 2n$ the bracket width contracts by $2^d/\|e\|_2^2 = d_{\mathrm{in}}(\square) = 4$, while the Verlet-parabola rational sine operator $\sin_{\mathrm{rat}}(d;\,p,q) = \frac{d_{\mathrm{in}}(\square)\,p(\|e\|_2^2 q - p)}{d_{\mathrm{out}}(E)(d)\,q^2 - p(\|e\|_2^2 q - p)}$ at $p/q\in\{1/d^2,\,1/d,\,1/\|e\|_2^2,\,\|e\|_2^2/d,\,1\}$ ($10^\circ,30^\circ,45^\circ,60^\circ,90^\circ$) returns $\frac{17}{97},\,\frac{1}{2},\,\frac{12}{17},\,\frac{32}{37},\,1$.
\item \textbf{Exact $O(|\mathbf{k}|^2)$ spherical isotropy, the $O(|\mathbf{k}|^4)$ fourth-moment anisotropy law, and the root-free 3D Euler-angle rotator $(\varphi_1,\varphi_2,\varphi_3)$ on the cuboctahedron:}
  for any wavevector $\mathbf{k}=(k_1,k_2,k_3)$,
  \begin{equation}
    \label{eq:cuboct-second-fourth-moments}
    \sum_{e\in O_2(d)} (\mathbf{k}\cdot e)^2 \;=\; 8\,|\mathbf{k}|^2,\qquad
    \sum_{e\in O_2(d)} (\mathbf{k}\cdot e)^4 \;=\; 12\,|\mathbf{k}|^4 - 4\sum_{i=1}^3 k_i^4,
  \end{equation}
  while composing three planar Weierstrass Pythagorean rotations $\mathcal{R}_{\mathrm{Euler}}(\varphi_1,\varphi_2,\varphi_3)$ across Euler angles $(\varphi_1,\varphi_2,\varphi_3)\in\mathbb{Z}_{N_{\mathrm{ring}}(d)}^3$ preserves the spherical norm $\|\mathcal{R}_{\mathrm{Euler}}(\varphi_1,\varphi_2,\varphi_3)\mathbf{v}\|_2^2 = \bigl(R_d(\varphi_1)R_d(\varphi_2)R_d(\varphi_3)\bigr)^2\|\mathbf{v}\|_2^2$ and tilts $(1,1,0)$ at $\varphi_2 = K_P(d)=64$ onto $(0,1,1)$.
\end{enumerate}
\end{theorem}

\section{Formal Verification in Lean~4}
\label{sec:lean4}

The formalization of Theorems~\ref{thm:main-1}--\ref{thm:main-3} (\texttt{thm1\_constant\_wronskian\_weierstrass\_lift\_and\_newton\_leibniz}, \texttt{thm2\_closed\_form\_fresnel\_jacobi\_dirichlet\_parseval\_bessel\_integrals}, and \texttt{thm3\_continuum\_integration\_from\_cuboctahedron\_archimedes\_layers}) resides in \texttt{RotatorCalculusPaper.lean} in \url{https://github.com/AI-Guiders/scientific-papers} (\texttt{0 sorryAx}, \texttt{0 Classical.choice}).

\begin{thebibliography}{9}
\bibitem{Malcev1973}
A.\,I.~Mal'cev, \emph{Algebraic Systems}, Die Grundlehren der mathematischen Wissenschaften, vol.~192, Springer-Verlag, Berlin--Heidelberg--New York, 1973 (Russian original: \emph{Алгебраические системы}, Nauka, Moscow, 1970).
\bibitem{Krause1986}
E.\,F.~Krause, \emph{Taxicab Geometry: An Adventure in Non-Euclidean Geometry}, Dover Publications, New York, 1986.
\bibitem{Thompson2011}
K.\,P.~Thompson, ``The nature of length, area, and volume in taxicab geometry,'' \emph{arXiv preprint} \texttt{arXiv:1101.2922} [math.MG], 2011.
\bibitem{Thompson2012}
K.\,P.~Thompson, ``Taxicab calculus: Trig derivatives,'' \emph{arXiv preprint} \texttt{arXiv:1205.3814v2} [math.MG], 2012.
\bibitem{Stroud1971}
A.\,H.~Stroud, \emph{Approximate Calculation of Multiple Integrals}, Prentice-Hall, Englewood Cliffs, NJ, 1971.
\bibitem{Moura2021}
L.~de~Moura and S.~Ullrich, ``The Lean~4 theorem prover and programming language,'' in \emph{CADE~28}, LNCS vol.~12699, Springer, 2021, pp.~625--635.
\end{thebibliography}

\end{document}
